\chapter{Sluttevaluering fra HDO}
\label{app:hdo-eval}

% Kvantitative verdier nedenfor er kilden til søylediagrammet i kap.~\ref{sec:evaluering}, figur~\ref{fig:hdo-eval-summary}.

Dette vedlegget gjengir den fullstendige sluttevalueringen fra
oppdragsgiver. Spørreskjemaet ble besvart av to representanter
fra HDO ved prosjektslutt. Kvantitative spørsmål er besvart på en skala fra 1 til 5, der 5
betyr fullstendig oppfylt. Antall respondenter er
begrenset, og evalueringen må derfor leses som en kvalitativ
ekspertvurdering fra oppdragsgiver, ikke som en statistisk
representativ undersøkelse.

\begin{table}[!htbp]
    \centering
    \footnotesize
    \renewcommand{\arraystretch}{1.25}
    \caption{Sluttevaluering fra HDO. Tallverdier er på skala
    1--5; tekstsvar er gjengitt ordrett. Tomme felter angir at
    respondenten ikke besvarte spørsmålet.}
    \label{tab:hdo-eval-full}
    \begin{tabular}{p{6.2cm} p{4.0cm} p{4.0cm}}
        \toprule
        \textbf{Spørsmål} & \textbf{Respondent 1} &
        \textbf{Respondent 2} \\
        \midrule
        \multicolumn{3}{l}{\textit{Kravoppfyllelse}} \\
        \midrule
        Dekning av funksjonelle krav (F1.1--F1.5) &
        5 & 5 \\
        Savnet funksjonalitet eller prioriteringer for videre
        utvikling &
        --- & --- \\
        Plugin-arkitektur: bytting mellom backends tilstrekkelig
        isolert &
        5 & 5 \\
        Autentiseringsløsning (OAuth2/JWT, romtoken) for
        AMK/LVS-kontekst (S1.1--S1.2) &
        4 & 5 \\
        Logging av drifts- og sikkerhetshendelser (O1.1)
        tilstrekkelig &
        5 & 5 \\
        Prometheus-metrikker og health-endepunkter for
        driftsapparatet &
        «Metrics og endepunkt er bra og nødvendig for driften.» &
        «Ja, fungerer som forventet.» \\
        API-dokumentasjon (OpenAPI/Swagger) for leverandør
        (D1.1) &
        4 & 5 \\
        Løsning klar for videre integrasjon i HDOs
        videoløsninger &
        «Ja, kan tas videre uten store endringer.» &
        «Ja, men krever noe videreutvikling.» \\
        \midrule
        \multicolumn{3}{l}{\textit{Samarbeid og prosess}} \\
        \midrule
        Kommunikasjon med prosjektgruppen gjennom perioden &
        5 & 5 \\
        Kravspesifikasjonen tydelig nok som grunnlag &
        «Håper den var tydelig, men er helt sikkert rom for
        forbedring der også.» &
        --- \\
        Forhold som burde vært avklart tidligere &
        --- & --- \\
        \midrule
        \multicolumn{3}{l}{\textit{Samlet vurdering}} \\
        \midrule
        Samlet: i hvilken grad innfrir leveransen
        forventningene &
        4 & 4 \\
        Det viktigste HDO tar med seg videre &
        «Kommer til å benytte kodebasen videre mot våre
        leverandører. Så det er nok det viktigste.» &
        «Kodebase, en god POC.» \\
        Andre kommentarer eller tilbakemeldinger &
        --- & --- \\
        \bottomrule
    \end{tabular}
\end{table}
